\section{Unrestricted learning at bounded widths}
\label{sec:finitewidth-unrestricted}

A robustness restriction is unnecessary when the learner's width is itself
polynomially bounded. The following elementary argument retains the original
all-input clean margin and Gaussian conditioning conventions, and is uniform
in the Gaussian scale. It also covers analog target networks at their original
depth, without digital simulation.

Let $P_N=N(d+2)+(L-2)N^2$ be the parameter count, and let $\mu_{N,t}$ denote
the raw Gaussian with variances $t^2/d,t^2/N,t^2/N^2$, respectively. Thus
$\kappa_N=(t^2/N)\chi^2_{P_N}$ under $\mu_{N,t}$. Fix integers
$1\le C\le M$ and suppose, for some $D\ge1$, that
\begin{equation}\label{eq:finitewidth-targetmass}
 \mu_{N,t}\{w:f(x)z_w(x)>1/2\text{ for every }x\}
 \ge \exp(-DN/t^2).
\end{equation}
This holds for $C\le N\le M$ and $0<t\le1$. Only this target-mass
hypothesis is needed; posterior networks may be fragile.

\begin{theorem}\label{thm:finitewidth-unrestricted}
For $0<\varepsilon,\delta<1$, set
\begin{align*}
 B&=4D+2[d+2+(L-2)M]\log2+4\log(2/\varepsilon),\\
 A&=\sqrt{MB},\qquad T=d+1+M(L-1),\qquad
 q=[4T(MA)^{L-1}]^{-1}.
\end{align*}
If
\[
 n\ge\frac2\varepsilon\left[
 \log M+P_M\log(5A/q)+\log(1/\delta)\right],
\]
then, with probability at least $1-\delta$ over independent examples from any
input distribution, the unrestricted Gaussian posterior has Gibbs error at
most $\varepsilon$, simultaneously for every $C\le N\le M$ and $0<t\le1$.
Here ``unrestricted'' retains the original conditioning on
$|z_w(x)|>1/2$ for all inputs, but imposes no internal robustness condition.
The same sample event works after any number of additional correctly labeled
examples. In particular, if $D$ and $M$ are polynomial in $d,L,C$, the sample
bound is polynomial in $d,L,C,\varepsilon^{-1},\log(1/\delta)$, uniformly for
all $\gamma\ge1$ with $t^2=N^{1-\gamma}$.
\end{theorem}

\begin{proof}
For every sample, the event used to condition the raw Gaussian contains the
target event in~\eqref{eq:finitewidth-targetmass}. The chi-square exponential
moment therefore gives
\[
 Q_{S,N,t}\{\kappa_N>B\}
 \le\exp\left(\frac{DN}{t^2}+\frac{P_N}2\log2-
                       \frac{NB}{4t^2}\right)
 \le\varepsilon/2.
\]
The last inequality uses $P_N/N\le d+2+(L-2)M$, $N\ge1$, and $t\le1$.
This bound holds for every sample, without a probabilistic qualification.

If $\kappa_N\le B$, every physical parameter has magnitude at most $A$.
Round each to the nearest multiple of $q$, using at most $5A/q$ grid values
per coordinate. For the maximum coordinatewise hidden-state discrepancies,
boundedness and the Lipschitz property of tanh give
\[
 e_1\le(d+1)q,\qquad e_s\le NAe_{s-1}+Nq,\qquad
 |z-z_{\rm grid}|\le NAe_{L-1}+Nq.
\]
Consequently, uniformly over Boolean inputs,
\[
 |z-z_{\rm grid}|\le q[d+1+M(L-1)](MA)^{L-1}=1/4.
\]
Rounding thus preserves the entire Boolean function of every posterior
network with energy at most $B$.

The union of these grid classifiers over $C\le N\le M$ has logarithmic
cardinality at most $\log M+P_M\log(5A/q)$. A consistency union bound gives
one sample event on which every consistent grid classifier has true error
at most $\varepsilon/2$. Every good-energy posterior draw has the same
function as such a classifier; the remaining posterior mass is at most
$\varepsilon/2$. This proves the Gibbs bound. The finite class is independent
of $t$, and consistency only becomes more restrictive as examples are added.
\end{proof}

The logarithm of the grid size satisfies
\[
 \log(5A/q)=\log(20AT)+(L-1)\log(MA).
\]
The exponentially small parameter precision costs only a polynomial in
$L$ in this logarithm. The unresolved unrestricted problem is therefore
uniformity at learner widths \emph{larger than every polynomial in the
specified target complexity and depth}; this theorem does not show those
larger widths actually require more samples.

\begin{corollary}[Same-depth bounded-weight analog targets]
Suppose a width-$C$, depth-$L$ tanh network realizes $f$ with
$f(x)z(x)\ge1$ on every Boolean input and weights of magnitude at most
$R\ge1$. Fix $M\ge C$ and set
\begin{align*}
 K&=2[d(d+1)+(L-1)C^2]R^2,\\
 a&=[4(d+1+M(L-1))(MR)^{L-1}]^{-1},\\
 D&=\max\{1,K+[d+2+(L-2)M][2+\log(1/a)]\}.
\end{align*}
Then~\eqref{eq:finitewidth-targetmass} holds with this $D$.
Thus Theorem~\ref{thm:finitewidth-unrestricted} supplies a polynomial sample
bound at the original depth for every $C\le N\le M$, with dependence on
$R$ only through $\log R$ in the final sample bound. No robustness condition
on the target implementation or prior support is needed. A target with the
original strict margin $>1/2$ satisfies the displayed hypothesis after
doubling its readout.
\end{corollary}

\begin{proof}
Clone the $C$ neurons into balanced nonempty groups of sizes $n_i$ summing
to $N$. Copy first-layer rows and replace internal and readout weights by
$W_{ij}/n_j$ and $v_i/n_i$. This preserves the exact function. Since
$\max_i n_i/\min_i n_i\le2$, its energy is at most $K$, and its weights
remain bounded by $R$. A physical-coordinate perturbation of magnitude
at most $a$ changes every output by at most $1/4$, by the recurrence above.

Use isotropic coordinates
$y=(\sqrt d\,W^1,\sqrt N\,W^{\rm internal},Nv)$, so
$y\sim\mathcal N(0,t^2I)$ and $\|y_*\|^2\le NK$ for the cloned center.
The standard-normal-coordinate box
$|y_j/t-y_{*,j}/t|\le a$ lies within the physical box, because all coordinate
scalings are at least one and $t\le1$. Its probability is at least
\begin{align*}
 &\left(\frac{2a}{\sqrt{2\pi}}\right)^{P_N}
 \exp\left[-\frac12\sum_j(|y_{*,j}|/t+a)^2\right]\\
 &\qquad\ge\exp[-NK/t^2-P_N(2+\log(1/a))]
 \ge\exp(-DN/t^2).
\end{align*}
Here $(u+a)^2\le2u^2+2a^2$ and $a\le1$ suffice. Every network in this box
has target margin at least $3/4$.
\end{proof}

This is an elementary finite-parameter learning result, rather than a new
claim about the general VC dimension of neural networks. Its role is to
isolate which width regime remains open under the exact prior conventions.
